Els sistemes planetaris tendeixen a formar-se en ressonància. Això vol dir que, quan un planeta completa $n$ òrbites, el planeta veí en completa $m$, en què $n$ i $m$ són nombres enters. El novembre de 2023 la revista \textit{Nature} publicava el descobriment de sis planetes que orbitaven en ressonància al voltant de l'estrella HD110067.

\begin{apartat}{1,25}
Suposem que aquesta estrella té la massa del Sol i que els planetes tenen òrbites circulars. El sisè planeta, el més exterior, té un període de 54,7 dies. Calculeu la distància entre el planeta i l'estrella. Representeu l'estrella i el planeta, dibuixeu el vector d'acceleració normal i calculeu-ne el mòdul.
\end{apartat}

\begin{apartat}{1,25}
El cinquè planeta completa 4 òrbites en el mateix temps que el sisè planeta en completa 3 (relació 4:3). Calculeu el radi de l'òrbita del cinquè planeta i la seva energia mecànica suposant que la seva massa és 2,5 vegades la terrestre.
\end{apartat}

\dades{%
  $G = 6,67\times 10^{-11}\ \mathrm{N\,m^2\,kg^{-2}}$.\\
  Massa de la Terra, $M_\mathrm{T} = 5,98\times 10^{24}\ \mathrm{kg}$.\\
  Massa del Sol, $M_\mathrm{S} = 1,98\times 10^{30}\ \mathrm{kg}$.}
